\section{Implementación de ReAct Agent en LangGraph}

\subsection{Arquitectura basada en grafos}

La implementación de \texttt{create\_react\_agent} de LangGraph es más moderna. Basta con definir el modelo y las herramientas:

\begin{importantbox}{Lo esencial para entender LangGraph: los bordes virtuales (bordes condicionales)}
El ReAct Agent de LangGraph contiene dos nodos centrales:
\begin{itemize}
    \item \textbf{Nodo Agent}: se encarga de producir las solicitudes de llamada a herramientas
    \item \textbf{Nodo Tools}: se encarga de ejecutar las herramientas y devolver los resultados
\end{itemize}
El nodo Agent tiene un \textbf{borde condicional} que decide si se detiene la llamada a herramientas y se entrega directamente la respuesta, o si se continúa llamando a herramientas.
\end{importantbox}

\subsection{Gestión basada en una lista de mensajes}

A diferencia de LangChain, LangGraph gestiona la interacción por completo mediante \textbf{mensajes de conversación de múltiples turnos}:

Flujo de mensajes: Human Message $\rightarrow$ AI Message (con tool\_calls) $\rightarrow$ Tool Message(s) $\rightarrow$ AI Message (respuesta final)

\begin{knowledgebox}{Llamadas paralelas a herramientas}
LangGraph admite \textbf{llamadas paralelas a herramientas}: el campo \texttt{tool\_calls} de un AI Message es una lista, por lo que se pueden llamar varias herramientas al mismo tiempo. Por ejemplo, al consultar «la ciudad natal del campeón del Abierto de Australia», el modelo podría llamar en paralelo dos veces a la herramienta de búsqueda.
\end{knowledgebox}

\subsection{Resumen de la sección}
La implementación de ReAct en LangGraph se basa en una estructura de grafo y una lista de mensajes de múltiples turnos, con una arquitectura más clara y compatibilidad con llamadas paralelas a herramientas.

